\begin{tomtat}
	\textbf{ĐỀ 01}
\end{tomtat}
\begin{ex}
	Rút gọn biểu thức $B = \left(\dfrac{\sqrt{x}}{\sqrt{x}-1}-\dfrac{1}{x-\sqrt{x}}\right):\left(\dfrac{1}{\sqrt{x}+1}+\dfrac{2}{x-1}\right)$ với $x>0$ và $x \ne 1$.
	\loigiai{
		Với điều kiện $x>0$ và $x \ne 1$, ta có:
		\begin{align*} B &= \left(\dfrac{\sqrt{x}}{\sqrt{x}-1}-\dfrac{1}{\sqrt{x}(\sqrt{x}-1)}\right):\left(\dfrac{1}{\sqrt{x}+1}+\dfrac{2}{(\sqrt{x}-1)(\sqrt{x}+1)}\right) \\ &= \left(\dfrac{\sqrt{x} \cdot \sqrt{x} - 1}{\sqrt{x}(\sqrt{x}-1)}\right):\left(\dfrac{\sqrt{x}-1+2}{(\sqrt{x}-1)(\sqrt{x}+1)}\right) \\ &= \left(\dfrac{x-1}{\sqrt{x}(\sqrt{x}-1)}\right):\left(\dfrac{\sqrt{x}+1}{(\sqrt{x}-1)(\sqrt{x}+1)}\right) \\ &= \left(\dfrac{(\sqrt{x}-1)(\sqrt{x}+1)}{\sqrt{x}(\sqrt{x}-1)}\right):\left(\dfrac{\sqrt{x}+1}{(\sqrt{x}-1)(\sqrt{x}+1)}\right)\end{align*}
		Do $x > 0$ và $x \ne 1$ nên $\sqrt{x}-1 \ne 0$ và $\sqrt{x}+1 \ne 0$. Ta rút gọn các biểu thức trong ngoặc:
		\begin{align*} B &= \left(\dfrac{\sqrt{x}+1}{\sqrt{x}}\right) : \left(\dfrac{1}{\sqrt{x}-1}\right) \\ &= \dfrac{\sqrt{x}+1}{\sqrt{x}} \cdot (\sqrt{x}-1) \\ &= \dfrac{(\sqrt{x}+1)(\sqrt{x}-1)}{\sqrt{x}} \\ &= \dfrac{x-1}{\sqrt{x}}\end{align*}
		Vậy, biểu thức rút gọn là \textbf{$B = \dfrac{x-1}{\sqrt{x}}$}.
	}
\end{ex}

\begin{ex}
	Biết phương trình $2x^2 + 4x - 6 = 0$ có hai nghiệm $x_1$ và $x_2$, không giải phương trình, tính giá trị của biểu thức $A = \dfrac{1}{x_1} + \dfrac{1}{x_2} - 4x_1 - 4x_2$.
	\loigiai{
		Phương trình đã cho là $2x^2 + 4x - 6 = 0$.
		Ta có $a=2$, $b=4$, $c=-6$.
		Vì $ac = 2 \cdot (-6) = -12 < 0$, phương trình có hai nghiệm phân biệt $x_1, x_2$.
		Do $c = -6 \ne 0$ nên hai nghiệm $x_1, x_2$ đều khác 0 (điều này đảm bảo các biểu thức $\dfrac{1}{x_1}$ và $\dfrac{1}{x_2}$ có nghĩa).
		
		Áp dụng hệ thức Vi-ét, ta có:
		$$ \begin{cases}
			x_1 + x_2 = -\dfrac{b}{a} = -\dfrac{4}{2} = -2 \\
			x_1 x_2 = \dfrac{c}{a} = \dfrac{-6}{2} = -3
		\end{cases} $$
		Ta biến đổi biểu thức A:
		\begin{align*} A &= \dfrac{1}{x_1} + \dfrac{1}{x_2} - 4x_1 - 4x_2 \\ &= \left(\dfrac{x_2 + x_1}{x_1x_2}\right) - 4(x_1+x_2) \\ &= \dfrac{x_1+x_2}{x_1x_2} - 4(x_1+x_2)\end{align*}
		Thay các giá trị $x_1+x_2 = -2$ và $x_1x_2 = -3$ vào biểu thức A:
		\begin{align*} 
			A &= \dfrac{-2}{-3} - 4(-2) \\ &= \dfrac{2}{3} + 8 \\ &= \dfrac{2}{3} + \dfrac{24}{3} \\ &= \dfrac{26}{3}
		\end{align*}
		Vậy, giá trị của biểu thức là \textbf{$A = \dfrac{26}{3}$}.
	}
\end{ex}

\begin{ex}
	Giải hệ phương trình: 
	$\begin{cases} 2x - y = 6 \\ 3x + y = 14 \end{cases}$.
	\loigiai{
		Ta có hệ phương trình:
		$$ \begin{cases} 2x - y = 6 \quad &(1) \\ 3x + y = 14 \quad &(2) \end{cases} $$
		Cộng vế theo vế của phương trình (1) và phương trình (2), ta được:
		\begin{align*} (2x - y) + (3x + y) &= 6 + 14 \\ 2x + 3x - y + y &= 20 \\ 5x &= 20 \\ x &= \frac{20}{5} \\ x &= 4 \end{align*}
		Thay $x=4$ vào phương trình (2) (hoặc phương trình (1)):
		\begin{align*} 3(4) + y &= 14 \\ 12 + y &= 14 \\ y &= 14 - 12 \\ y &= 2 \end{align*}
		Vậy, nghiệm của hệ phương trình là $(x; y) = (4; 2)$.
		Hoặc viết:
		$$ \begin{cases} x = 4 \\ y = 2 \end{cases} $$
	}
\end{ex}
\begin{ex}
	Viết ngẫu nhiên một số tự nhiên có hai chữ số nhỏ hơn hoặc bằng 84 và lớn hơn 60. Tính xác suất số tự nhiên viết ra là số chia cho 3 và chia cho 4 đều dư 1.
	\loigiai{
		\begin{enumerate}
			\item Tập hợp tất cả các số tự nhiên có hai chữ số nhỏ hơn hoặc bằng 84 và lớn hơn 60 (kết quả có thể xảy ra) là: \\ $\{61, 62, 63, 64, 65, 66, 67, 68, 69, 70, 71, 72, 73, 74, 75, 76, 77, 78, 79, 80, 81, 82, 83, 84\}$. \\
			Số kết quả có thể xảy ra là 24.
			\item Các số trong tập hợp trên mà chia cho 3 và chia cho 4 đều dư 1 (kết quả thuận lợi cho biến cố) là: \\ $\{61, 73\}$. \\
			Số kết quả thuận lợi cho biến cố là 2.
			\item Xác suất của biến cố là: \\
			$P = \dfrac{\text{Số kết quả thuận lợi}}{\text{Số kết quả có thể xảy ra}} = \dfrac{2}{24} = \dfrac{1}{12}$.
		\end{enumerate}
		Vậy, xác suất cần tính là \textbf{$\dfrac{1}{12}$}.
	}
\end{ex}

\begin{ex}
	Một tổ sản xuất phải làm xong 4800 sản phẩm trong một số ngày quy định. Thực tế, mỗi ngày tổ đó làm được nhiều hơn 100 sản phẩm theo kế hoạch. Vì thế 8 ngày trước khi hết hạn, tổ sản xuất đã làm xong 4800 sản phẩm. Hỏi theo kế hoạch, mỗi ngày tổ sản xuất làm bao nhiêu sản phẩm? (Biết rằng mỗi ngày tổ sản xuất làm số sản phẩm là như nhau).
	\loigiai{
		Gọi số sản phẩm tổ sản xuất làm mỗi ngày theo kế hoạch là $x$ (sản phẩm).
		Điều kiện: $x > 0$.
		
		Thời gian dự định để hoàn thành 4800 sản phẩm theo kế hoạch là $\dfrac{4800}{x}$ (ngày).
		
		Thực tế, mỗi ngày tổ sản xuất làm được $x + 100$ (sản phẩm).
		Thời gian thực tế để hoàn thành 4800 sản phẩm là $\dfrac{4800}{x+100}$ (ngày).
		
		Theo đề bài, tổ sản xuất đã làm xong 4800 sản phẩm sớm hơn 8 ngày so với kế hoạch, nên ta có phương trình:
		$$ \dfrac{4800}{x} - \dfrac{4800}{x+100} = 8 $$
		$$ x^2 + 100x - 60000 = 0 $$
		Giải phương trình bậc hai này, ta tìm được hai nghiệm:
		$x_1 = 200$
		$x_2 = -300$
		
		Đối chiếu với điều kiện $x > 0$:
		\begin{itemize}
			\item $x = 200$ (thỏa mãn điều kiện $x > 0$)
			\item $x = -300$ (không thỏa mãn điều kiện $x > 0$, nên loại)
		\end{itemize}
		Vậy, theo kế hoạch, mỗi ngày tổ sản xuất làm 200 sản phẩm.
	}
\end{ex}

\begin{ex}
	Cho tam giác $ABC$ vuông tại $A$. Vẽ đường tròn tâm $C$, bán kính $CA$. Từ điểm $B$ kẻ tiếp tuyến $BM$ với đường tròn $(C; CA)$ ($M$ là tiếp điểm, $M$ và $A$ nằm khác phía đối với đường thẳng $BC$).
	\begin{enumerate}
		\item[a)] Chứng minh tứ giác ACMB nội tiếp.
		%\item[b)] Lấy điểm $N$ thuộc đoạn thẳng $AB$ ($N \neq A, N \neq B$). Lấy điểm $P$ thuộc tia đối của tia $MB$ sao cho $MP = AN$. Chứng minh tam giác $CPN$ là tam giác cân. Đường thẳng $AM$ đi qua trung điểm của đoạn thẳng $NP$.
	\end{enumerate}
	\loigiai{
		\begin{center}
			\includegraphics[scale=0.1]{hinh/1}
		\end{center}
		* Xét $\triangle CAN$ và $\triangle CMP$ có:
		
		$CA = CM$; $\widehat{CAN} = \widehat{CMP} = 90^\circ$; $AN = MP$
		
		$\Rightarrow \triangle CAN \sim \triangle CMP$ (c.g.c) $\Rightarrow CN = CP$
		
		$\Rightarrow$ Tam giác CPN cân tại C
		
		* Gọi I là trung điểm của đoạn thẳng NP
		
		Tam giác CPN cân tại C và I là trung điểm của đoạn thẳng NP nên $CI \perp NP$
		
		Tứ giác NACI nội tiếp $\Rightarrow \widehat{NIA} = \widehat{NCA}$
		
		Tứ giác CIMP nội tiếp $\Rightarrow \widehat{MIP} = \widehat{MCP}$
		
		$\triangle CAN \sim \triangle CMP \Rightarrow \widehat{NCA} = \widehat{MCP}$
		
		Ta có $\widehat{NIA} + \widehat{PIA} = 180^\circ$ (vì I nằm giữa N và P)
		
		$\Rightarrow \widehat{MIP} + \widehat{PIA} = 180^\circ$ mà 2 góc này kề nhau
		
		$\Rightarrow$ A, I, M là 3 điểm thẳng hàng
		
		KL: Đường thẳng AM đi qua trung điểm của đoạn thẳng NP
		}
\end{ex}
